% Generator: Claude Opus 5 (High)
% Glossary of the English/German term pairs mirrored on first introduction throughout the
% document, as required by the "German mirroring" rule in gemini.md. Every entry below
% corresponds to exactly one \newterm{...} (\germanterm{...}) pair in content/week-NN.tex,
% and the "Week" column records where that first introduction happens.
%
% MAINTENANCE: when adding a \germanterm{...} anywhere, add its row here too. To list every
% pair currently in the document and compare against this table:
%   grep -ohE '\\germanterm\{[^}]*\}' content/week-*.tex | sort -u
% A German term must appear exactly once across all week files (first introduction only);
% duplicates are a style error, not a harmless repetition.

\chapter{Glossary --- English and German terms}
\label{ch:glossary}

The lecture and the exercise sheets for this course exist in both English and German, and the
two vocabularies do not always line up in the obvious way --- \germanterm{Abschluss} is the
closure and not a \qt{conclusion}, \germanterm{Rand} is the boundary and not a \qt{rim}, and
\germanterm{Satz von Schwarz} is what English-language texts usually call Clairaut's theorem.
Each term below is mirrored once in the main text, at the point where it is first introduced;
this appendix collects those pairs in one place so the document can be used alongside German
material.

Canonical German wording follows Jérôme Paschoud's topic-named files. The \textbf{Week} column
gives the chapter of first introduction.

\section*{A --- E}

\begin{center}
\begin{tabular}{@{}p{0.40\linewidth} p{0.44\linewidth} c@{}}
\textbf{English} & \textbf{German} & \textbf{Week} \\
\hline
accumulation point & Häufungspunkt & 2 \\
antisymmetric $k$-linear form & antisymmetrische $k$-lineare Form & 11 \\
boundary & Rand & 2 \\
bounded & beschränkt & 3 \\
chain rule & Kettenregel & 4 \\
closed & abgeschlossen & 2 \\
closure & Abschluss & 2 \\
compact & kompakt & 2 \\
complete metric space & vollständiger metrischer Raum & 2 \\
connected & zusammenhängend & 3 \\
continuity equation & Kontinuitätsgleichung & 10 \\
continuous & stetig & 2 \\
contraction & Kontraktion & 3 \\
convex & konvex & 6 \\
critical point & kritischer Punkt & 5 \\
cuspidal cubic & Neilsche Parabel & 5 \\
degenerate (critical point) & entartet & 5 \\
diffeomorphism & Diffeomorphismus & 6 \\
differentiable & differenzierbar & 4 \\
differential (the) & das Differential & 4 \\
differential $p$-form & Differentialform & 11 \\
directional derivative & Richtungsableitung & 4 \\
disconnected & unzusammenhängend & 3 \\
\end{tabular}
\end{center}

\section*{F --- L}

\begin{center}
\begin{tabular}{@{}p{0.40\linewidth} p{0.44\linewidth} c@{}}
\textbf{English} & \textbf{German} & \textbf{Week} \\
\hline
figure eight & Lemniskate & 8 \\
fixed point & Fixpunkt & 3 \\
flux & Fluss & 10 \\
Frobenius norm & Frobeniusnorm & 2 \\
Gabriel's horn & Torricellis Trompete & 10 \\
Gram determinant & Gram-Determinante & 9 \\
gradient & Gradient & 4 \\
Hessian matrix & Hesse-Matrix & 4 \\
Hölder's inequality & Höldersche Ungleichung & 1 \\
homeomorphism & Homöomorphismus & 6 \\
Hurwitz criterion & Hurwitz-Kriterium & 5 \\
isometry & Isometrie & 9 \\
Jacobi matrix & Jacobi-Matrix & 4 \\
Jordan measurable & Jordan-messbar & 8 \\
Lagrangian & Lagrange-Funktion & 5 \\
Lebesgue number & Lebesgue-Zahl & 2 \\
Lebesgue's covering lemma & Lebesguesches Überdeckungslemma & 2 \\
length & Länge & 9 \\
level set & Niveaumenge & 5 \\
Lipschitz continuous & Lipschitz-stetig & 3 \\
local minimum & lokales Minimum & 5 \\
\end{tabular}
\end{center}

\section*{M --- P}

\begin{center}
\begin{tabular}{@{}p{0.40\linewidth} p{0.44\linewidth} c@{}}
\textbf{English} & \textbf{German} & \textbf{Week} \\
\hline
metric & Metrik & 2 \\
metric space & metrischer Raum & 2 \\
multi-index & Multiindex & 4 \\
negative definite & negativ definit & 5 \\
negative semidefinite & negativ semidefinit & 6 \\
norm & Norm & 2 \\
normal space & Normalraum & 7 \\
open & offen & 2 \\
open cover & offene Überdeckung & 2 \\
orientable & orientierbar & 12 \\
partial derivative & partielle Ableitung & 4 \\
path & Weg & 3 \\
path-connected & wegzusammenhängend & 3 \\
positive definite & positiv definit & 5 \\
positive semidefinite & positiv semidefinit & 6 \\
pseudo-metric & Pseudometrik & 2 \\
pullback & Zurückziehung & 12 \\
\end{tabular}
\end{center}

\section*{Q --- Z}

\begin{center}
\begin{tabular}{@{}p{0.40\linewidth} p{0.44\linewidth} c@{}}
\textbf{English} & \textbf{German} & \textbf{Week} \\
\hline
saddle point & Sattelpunkt & 5 \\
scalar field & Skalarfeld & 4 \\
inner product / scalar product & Skalarprodukt & 2 \\
Schwarz's theorem & Satz von Schwarz & 5 \\
sequence & Folge & 1 \\
shadow price & Schattenpreis & 5 \\
space & Raum & 2 \\
submanifold & Untermannigfaltigkeit & 7 \\
submultiplicative & submultiplikativ & 3 \\
support (of a partition function) & Träger & 12 \\
tangent space & Tangentialraum & 7 \\
totally bounded & total beschränkt & 2 \\
trapezoidal rule (Gauss) & Gaußsche Trapezformel & 13 \\
uniform convergence & gleichmässige Konvergenz & 2 \\
uniformly continuous & gleichmässig stetig & 2 \\
velocity (of a curve) & Geschwindigkeit & 4 \\
wedge product & Dachprodukt & 11 \\
interior & Inneres & 2 \\
\end{tabular}
\end{center}

\begin{ainote}
Three conventions in this table are worth flagging, because they are places where the German
and English literature genuinely diverge rather than merely translate.
\begin{itemize}
  \item \textbf{\germanterm{Satz von Schwarz}} is the symmetry of mixed partial derivatives
    (\cref{thm:schwarz_mixed_partials}). English texts more often call it \emph{Clairaut's
    theorem}; both names are used in this document, with Schwarz first because that is what the
    lecture uses.
  \item \textbf{\germanterm{Neilsche Parabel}} names the curve $y^2 = x^3$ after William Neile.
    The English \emph{cuspidal cubic} is descriptive rather than eponymous, so searching for one
    name will not find the other.
  \item \textbf{\germanterm{gleichmässig}} carries both English senses --- it is
    \qt{uniform} in \emph{uniform convergence} and in \emph{uniformly continuous}, which are
    different notions (\cref{ex:uniform_convergence}, \cref{def:uniformly_continuous}). The
    German gives no hint that these are distinct, and neither does the English.
\end{itemize}
\end{ainote}
